% =====================================================================
%  §4 — Worked Examples from the Literature
%
%  Six-paper corpus table, two featured examples (Duerr 2006, Inouye
%  2018) as short prose, with full annotations hosted as case reports
%  in the project repository.
%
%  Revised 2026-04-21: YAML listings moved to repository case-reports
%  to save space; per-example prose retained the methodological points
%  that the listings previously anchored.
% =====================================================================

\section{Illustrative Examples from the Literature}
\label{sec:examples}

To demonstrate the model's coverage and identify where it reaches its
current limits, we applied it to six published genetic-evidence papers that
span the major epistemic shapes encountered in literature-based variant
interpretation (Supplementary Table~\ref*{supp:tab:corpus}). The set covers
classical Mendelian human genetics (Gupta), population-scale association
(Duerr), case--control resequencing (Davis), model-organism and \emph{in
vitro} functional genetics (Jossin, Nelson, and the functional components of
Davis), and a genome-wide polygenic score (Inouye). Four of the six papers
were annotated manually from curator PDF highlights and callouts and form
the curator-authored reference set for this feasibility corpus. The remaining two were drafted by an AI assistant
and reviewed by a curator. All annotations and per-paper case reports are
available in the project
repository.\footnote{\url{https://github.com/ForomePlatform/genetic-evidence-model/tree/master/case-reports}}

Duerr et al.\ 2006 (\cref{sec:ex-duerr}) is the reference example of a
classical variant-level association study. Inouye et al.\ 2018
(\cref{sec:ex-inouye}) is a polygenic risk score study that exposes
limitations of the current variant-level model and motivates a cluster of
candidate extensions (\cref{sec:discussion}).

\subsection{A Genome-Wide Association Study Example}
\label{sec:ex-duerr}

Duerr et al.\ 2006 identified \textit{IL23R} as an inflammatory bowel
disease gene through a genome-wide association study of non-Jewish patients
with ileal Crohn's disease, followed by replication in a Jewish cohort and
family-based transmission testing. Under our decomposition, the paper yields
seven \GeneticEvidence{} items (Supplementary
Table~\ref*{supp:tab:corpus} and Case Report CR-DU): variant-level discovery
and replication, a family-based transmission analysis, conditional
fine-mapping, negative results at related pathway genes, and cited
functional and mouse-model support.

Six of the seven items fit the current schema without strain. The seventh, a
cited anti-p40 antibody trial, concerns drug-response evidence for which the
model has no fitting genetic Knowledge Domain, revealing a
translational-scope gap. Review of this annotation also identified three
candidate extensions that remain provisional under the two-papers rule: an
explicit direction-of-effect dimension for protective and risk-conferring
alleles, a construct for cross-item therapeutic synthesis, and an extension
of conditional activation from gene-level to variant-level Mendelian context.
Further detail is given in Supplementary
Note~\ref*{supp:sec:supp-ex-duerr}.

\subsection{A Polygenic Risk Score Example}
\label{sec:ex-inouye}

Inouye et al.\ 2018 developed and externally validated a meta-genomic risk
score (metaGRS) for coronary artery disease by combining three component
scores and evaluating performance in UK Biobank participants. Unlike Duerr's
variant-level association study, the evidential target is a derived,
whole-genome score and the primary measurements are epidemiological.

Our annotation decomposes the paper into four \GeneticEvidence{} items:
predictive and stratifying performance; independence from conventional risk
factors; persistence of stratification among medicated individuals; and the
headline finding that combining three component scores outperforms each
component (Case Report CR-IN). All four require deliberate placeholder
assignments in the current model. In particular, the score is assigned the
placeholder Target Type \dimval{Variant}, and Variant Ascertainment is marked
not applicable. These forced fits are recorded as six candidate extensions:
a score target type, whole-genome-aggregate resolution, a target-composition
axis, epidemiological measurement targets, a structured cohort descriptor,
and a flag for derived-artifact targets. The model's limitations are thus
recorded explicitly rather than left implicit. Further detail is given in
Supplementary Note~\ref*{supp:sec:supp-ex-inouye}.
